\chapter{The One Theorem}
% OUTLINE Ch5 "The One Theorem (Lawvere 1969)". Laws: both inherited
% (labeled intuitions adds-no-claim); non-funge. CITATION: Lawvere 1969
% ("Diagonal arguments and cartesian closed categories") -- the
% fixed-point theorem from which Cantor/Godel/Tarski/Russell fall as
% instances; the family-resemblance ("roughly the same size") + the
% "always shows up" clauses EARNED here as the theorem's content. THE
% KEY DIFFERENCE FROM VOL 3's Ch5, stated in-chapter: there the unity was
% a READING; here it is a PROVED, CITED theorem -- and the chapter says
% exactly why it can claim more. Zero device appearances (dose = Ch6).
% Gate: Kodo (the Lawvere earn + the stronger-claim justification).

\begin{intuition}
Four times now the reader has watched the same trick. A list is turned
against itself, a diagonal is walked, a fixed point is forced --- Cantor's
missing number, Russell's rogue class, G\"odel's unprovable sentence,
Tarski's unnameable truth. The costumes differ; the choreography does not.
By the fourth performance the honest reaction is not admiration but
suspicion: these are not four theorems that happen to rhyme. They are one
theorem, and someone must have written it down. The feeling to carry:
when a gesture recurs this exactly, look for the single result of which
all the sightings are instances --- and this chapter is the one place in
the book where that suspicion is not merely voiced but discharged.
\end{intuition}

A word of contrast belongs here, because the previous volume ended on a
chapter shaped like this one and this chapter must be honest about the
difference. In \emph{The Intuition}, four faces of mathematics were shown
to be four forms of one gesture --- and that unity was fenced, carefully,
as a \emph{reading}: this author's arrangement of cited material, no
theorem of metaphysics, the object an image and the citations the content.
The reader was told, in so many words, that the oneness was a way of
seeing and not a proved fact. This chapter claims more, and can, because
the situation is genuinely different: the unification of the diagonal
arguments is not a reading a book imposes. It is a theorem someone proved,
with a date and a name, and this book's stronger claim here is not a
loosening of discipline but a reporting of the record. When the unity is
an author's arrangement, the honest word is ``reading.'' When the unity is
a published theorem, the honest word is ``theorem'' --- and saying so is
the same discipline, not a departure from it.

The name is F.\,William Lawvere, and the paper is his 1969 ``Diagonal
arguments and cartesian closed categories.'' Its achievement, stated
without its machinery, is this: Lawvere identified the abstract structure
common to every diagonal argument --- a very general condition on a
mapping under which a fixed point is forced --- and showed that Cantor's
theorem, G\"odel's incompleteness, Tarski's undefinability, Russell's
paradox, and the halting problem's unsolvability all fall out as instances
of a single fixed-point statement. Where the earlier chapters each built
the room by hand in one house, Lawvere exhibited the floor plan that every
such house is built to --- and proved that any house of the right kind
must contain the room. The diagonal is not a family of clever tricks that
happen to resemble one another; it is one theorem with many
interpretations, and Lawvere is who established that, in the register where
establishing means proving.

Two of the preface's four clauses are paid precisely here, and paid as the
theorem's own content rather than as observations across the chapters.
\emph{The room always shows up}: Lawvere's condition is mild --- it asks
only that a system be able to represent its own maps richly enough --- and
wherever it holds, the fixed point is not optional but forced. ``Always''
stops being an induction over four examples and becomes the conclusion of
one theorem: in any structure meeting the hypothesis, the room exists,
necessarily. And \emph{the room is, roughly, always the same size}: what
the four rooms literally share, Lawvere's proof makes exact --- the same
diagonal core, the same fixed-point-forcing step, dressed in each house's
particular furniture. The ``roughly'' is now locatable: the invariant part
is the fixed-point construction the theorem isolates; the varying part is
the interpretation each house supplies. Family resemblance, given a
precise family and a precise resemblance.

\begin{intuition}
The picture, and it is the previous volume's tesseract grown a dimension
more honest. There, four faces were four shadows of one object, and the
object was offered as an image the reader might hold. Here the object is
not an image; it is a theorem, and the four rooms are its four models the
way four numbers can all satisfy one equation. Same shadow-and-object
feeling --- but this time the object was caught, named, and published, and
the reader may check the citation rather than trust the picture. The
image became a reference.
\end{intuition}

The non-funge law survives this unification intact, and stating how is the
chapter's last and most careful step, because a careless reader will hear
``one theorem'' and conclude ``one room after all.'' It does not follow,
and the distinction is the hinge the next chapter turns on. Lawvere
unified the \emph{construction}: there is one floor plan, one theorem,
provably shared. He did not identify the \emph{rooms}: G\"odel's actual
undecidable sentence, the specific string built in the specific system,
remains that system's own, and Tarski's actual undefinable predicate
remains its language's own. One theorem generates a room in each house it
visits, and the rooms it generates are as distinct as the houses ---
built to one plan, standing in different places, furnished by different
occupants, and never the same room. The unity is in the plan; the
non-funge is in the buildings. A shared blueprint is exactly what lets a
thousand distinct houses each have their own front room, and it is the
opposite of a claim that they share one.

So the book arrives, with a theorem in hand, at the seam it has been
walking toward: the construction is one and provably so; the rooms are
many and bound each to its own argument. Hold both --- they are not in
tension, they are the same fact stated from two ends --- and the next
chapter can do the thing this whole volume was for: set the rooms side by
side, show that no door connects them, and state, as law, that you cannot
carry one across to another. One plan; no shared room. That is the sixth
chapter, and the place the series' own instance finally, and only, appears.
